\section*{Anotācija}
\addcontentsline{toc}{section}{Anotācija (LV)}

\textbf{Nosaukums:} Telegram Bot informācijas sistēmas izstrāde
sabiedriskajam transportam Fergānas pilsētā.

\medskip

Sabiedriskajam transportam Fergānā -- aptuveni 300\,000 iedzīvotāju
lielā pilsētā Uzbekistānā -- nav oficiālas digitālās informācijas
sistēmas. Iedzīvotāji un viesi paļaujas uz mutisku informāciju, papīra
zīmēm vai vispārīgiem karšu pakalpojumiem, piemēram, Yandex Maps, kas
nesniedz pilnīgus datus par vietējiem autobusu maršrutiem. Tas rada
ikdienas problēmas: cilvēki nokavē autobusus, izvēlas nepareizus
maršrutus vai pavada lieku laiku pieturās.

Darba mērķis ir izstrādāt Telegram bota informācijas sistēmu, kas
vienkāršā un pieejamā veidā nodrošina sabiedriskā transporta
informāciju Fergānas pilsētā. Telegram tika izvēlēts tāpēc, ka tas ir
viens no populārākajiem ziņojumapmaiņas lietotnēm Uzbekistānā un labi
darbojas pat uz lētiem viedtālruņiem ar ierobežotu internetu.

Bots tika izstrādāts valodā Python, izmantojot \texttt{aiogram}
ietvaru un SQLite datubāzi. Maršrutu un pieturu dati tika apkopoti no
Yandex Maps un saglabāti vietējā datubāzē. Bots atbalsta sešus
autobusu maršrutus un piedāvā šādas galvenās funkcijas: maršruta
pieturu rādīšanu, maršruta meklēšanu starp divām pieturām, taksometru
informāciju un lietotāju ieteikumu vākšanu. Sistēma tika izvietota
PythonAnywhere mākoņa platformā nepārtrauktai darbībai.

Bots tika testēts ar nelielu lietotāju grupu Fergānā. Rezultāti rāda,
ka sistēma ir viegli lietojama, ātri reaģē un aptver biežākos
transporta jautājumus. Darbā arī tiek aplūkoti turpmākie uzlabojumi,
tostarp integrācija ar Grok AI modeli dabiskās valodas vaicājumiem.

Darbs apstiprina, ka Telegram bots ir praktisks un zemu izmaksu
risinājums pilsētām, kurās nav oficiālu transporta lietotņu.

\medskip

\noindent\textbf{Apjoms:} XX lappuses, YY tabulas, ZZ attēli,
zx atsauces, xz pielikumi.
\newpage
